% ======================================================================
%  KAVRAM: Fine-tuning  (Tier 1)  — kompakt, tablo ağırlıklı
% ======================================================================

\concept{Training · Eğitim}{Fine-tuning}{Davranışı her çağrıda anlatmak yerine modelin ağırlıklarına işlemek}

\stage{1}{Problem}
Model istediğin gibi davranmıyor: format tutmuyor, marka tonunu yakalamıyor,
alan jargonunu bilmiyor. Çözüm diye her prompt'a aynı uzun talimatları, aynı 10
örneği dolduruyorsun --- token şişiyor, tutarlılık yine de oynuyor. Prompt
mühendisliği bir tavana çarpıyor.

\keyline{Soru: Bu davranışı her çağrıda yeniden \emph{anlatmak} mı, yoksa modele bir kez \emph{öğretmek} mi daha doğru?}

\stage{2}{Neden Var}
Bazı davranışlar prompt'ta tekrar tekrar açıklanmak yerine modelin kendisine
gömülmelidir. Fine-tuning, modeli kendi örneklerinle yeniden eğiterek
ağırlıklarını --- yani \emph{varsayılan davranışını} --- değiştirir. Böylece
talimat prompt'tan çıkar, modelin doğasına geçer.

\stage{3}{Tarihçe}
\begin{timeline}
\yr{2010'lar} tam yeniden eğitim / transfer learning \tlsep \yr{2021} LoRA / PEFT: ucuz, ``adapter'' tabanlı fine-tune \tlsep \yr{2022} instruction tuning + RLHF: model talimat izlemeyi öğrenir \tlsep \yr{2024–26} yönetilen fine-tune API'leri yaygınlaşır
\end{timeline}

\stage{4}{Tanım}
\begin{problembox}
\textbf{Fine-tuning:} Önceden eğitilmiş bir modeli, kendi örnek veri kümen
üzerinde eğitmeye devam ederek ağırlıklarını belirli bir davranış / üslup /
format için ayarlama. \textbf{LoRA / PEFT:} tüm modeli değil, küçük bir
``adapter'' katmanını eğiten ucuz varyant.
\end{problembox}

\stage{5}{Zihinsel Model}
\begin{mentalbox}
Prompt vermek, bir çalışana \textbf{her sabah talimat} vermektir. Fine-tuning,
onu bir \textbf{kursa göndermektir}: davranış artık ikinci doğası olur, her gün
anlatman gerekmez. Ama dezavantajı da kurs gibidir --- yeni \emph{bilgi} için
her seferinde yeniden kursa göndermek pahalıdır ve çalışan bu sırada genel
becerilerini kısmen ``unutabilir'' (\en{catastrophic forgetting}).
\end{mentalbox}

\stage{6}{Mimari Diyagram}
\vspace{2pt}
\begin{center}
\begin{tikzpicture}[font=\sffamily\scriptsize, node distance=7mm,
   box/.style={draw=accent, line width=0.8pt, fill=wash, rounded corners=2pt,
               minimum height=10mm, minimum width=22mm, align=center},
   io/.style={draw=softgray, fill=white, rounded corners=8pt, align=center,
              minimum height=10mm, minimum width=22mm, inner sep=3pt},
   flow/.style={-{Stealth[length=2mm]}, accent, line width=0.8pt}]
  \node[io] (base) {Temel model\\\textcolor{softgray}{genel}};
  \node[io, below=8mm of base] (data) {Örnek veri\\\textcolor{softgray}{girdi→çıktı çiftleri}};
  \node[box, right=12mm of base] (train) {Eğitim\\(LoRA / PEFT)};
  \node[draw=accentdk, fill=accent!12, rounded corners=2pt, right=12mm of train,
        minimum height=10mm, minimum width=24mm, align=center] (ft)
        {Fine-tuned\\model};
  \draw[flow] (base) -- (train);
  \draw[flow] (data.east) -- ++(6mm,0) |- (train.south west);
  \draw[flow] (train) -- (ft);
\end{tikzpicture}
\end{center}
\keyline{Dikkat: çıktı yeni bir model (veya adapter) ağırlığıdır --- bir indeks değil. Bilgi değişince yeniden eğitmek gerekir.}

\stage{7}{Karar Çerçevesi}
İhtiyacın \emph{bilgi} mi, \emph{davranış} mı? Tek ayrım kararın çoğunu verir.
\vspace{2pt}
\begin{center}
\begin{tikzpicture}[font=\sffamily\scriptsize, node distance=5mm and 8mm,
   q/.style={draw=accent, line width=0.8pt, fill=wash, diamond, aspect=2.4,
             align=center, inner sep=1pt, minimum width=30mm},
   leaf/.style={draw=accentdk, line width=0.8pt, rounded corners=2pt,
                fill=accent!12, align=center, inner sep=4pt, minimum width=22mm},
   yes/.style={good, font=\sffamily\bfseries\scriptsize},
   no/.style={warn, font=\sffamily\bfseries\scriptsize},
   ed/.style={-{Stealth[length=2mm]}, softgray, line width=0.7pt}]
  \node[q] (q1) {Değişen/güncel\\bilgi mi gerek?};
  \node[leaf, below left=13mm and -2mm of q1] (rag) {\textbf{RAG}\\(fine-tune değil)};
  \node[q, below right=13mm and -2mm of q1] (q2) {Yeterli kaliteli\\örnek var mı?};
  \node[leaf, below left=13mm and 4mm of q2, minimum width=20mm] (prompt) {\textbf{Önce prompt}\\/ few-shot};
  \node[leaf, below right=13mm and 4mm of q2, minimum width=20mm] (ft) {\textbf{Fine-tune}\\stil/format/davranış};
  \draw[ed] (q1) -- node[yes,left]{Evet} (rag);
  \draw[ed] (q1) -- node[no,right]{Hayır} (q2);
  \draw[ed] (q2) -- node[no,left]{Hayır} (prompt);
  \draw[ed] (q2) -- node[yes,right]{Evet} (ft);
\end{tikzpicture}
\end{center}
\keyline{Kural: Önce prompt'u tüket. Fine-tuning'i yalnızca tutarlı davranış/üslup/format için, kanıtlanmış ihtiyaç varken kullan.}

\stage{8}{Üretim Senaryosu}
\begin{prodbox}[Üretim Senaryosu · Katı Format + Marka Tonu]
\textbf{Bağlam.} Bir ürün, çıktıların \emph{her zaman} belirli bir JSON şemasında
ve sabit bir marka tonunda olmasını istiyor. Prompt + few-shot ile \%95'e
ulaşılıyor ama kalan \%5 hata production'da pahalı.

\smallskip
\textbf{Çözüm.} Birkaç bin temiz ``girdi→ideal çıktı'' çifti ile LoRA fine-tune.
Bilgi hâlâ RAG'den geliyor; fine-tune yalnızca \emph{biçim ve tonu} garantiliyor.

\smallskip
\textbf{Sonuç.} Format hataları neredeyse sıfırlanıyor, prompt kısalıyor (talimat
artık modelde), token maliyeti düşüyor. Bilgi RAG'de kaldığı için güncel.
\end{prodbox}

\stage{9}{Anti-pattern}
\begin{antibox}[Anti-pattern · ``Bilgi için Fine-tune'']
Şirket dokümanlarını modele ``öğretmek'' için fine-tune yapmak. Bilgi
ağırlıklara gömülünce: güncellenemez (her değişiklik = yeniden eğitim), kaynak
gösterilemez, ve model yine de emin emin halüsinasyon yapar.

\smallskip
\textbf{Belirti:} ``Fine-tune ettik ama cevaplar eski/yanlış ve kaynaksız.''\\
\textbf{Kök neden:} Bilgi problemi, davranış aracıyla çözülmeye çalışılmış.\\
\textbf{Düzeltme:} Bilgi $\to$ RAG; davranış/format/üslup $\to$ fine-tune. İkisini karıştırma.
\end{antibox}

\stage{10}{Alternatifler \& Ödünleşimler}
Bu, kitabın en sık sorulan kararıdır. Tek tabloda:\par
\vspace{2pt}
\noindent\begin{tabularx}{\linewidth}{@{}>{\raggedright\arraybackslash}p{24mm} L L L@{}}
\arrayrulecolor{hairline}\toprule
{\sffamily\bfseries\small } &
{\sffamily\bfseries\small Prompt / few-shot} &
{\sffamily\bfseries\small RAG} &
{\sffamily\bfseries\small Fine-tuning} \\
\midrule
\textbf{Çözdüğü} & Hızlı yönlendirme & Güncel \emph{bilgi} & Kalıcı \emph{davranış} \\
\addlinespace[2pt]
\textbf{Güncelleme} & Anında & Anında (veri) & Yeniden eğitim \\
\addlinespace[2pt]
\textbf{Maliyet} & Düşük & Orta (altyapı) & Yüksek (eğitim) \\
\addlinespace[2pt]
\textbf{Kaynak gösterir} & Hayır & Evet & Hayır \\
\addlinespace[2pt]
\textbf{En iyi} & İlk deneme & Bilgi yoğun & Stil/format/ton \\
\bottomrule
\end{tabularx}
\vspace{3pt}
Gerçek sistemler genelde \textbf{RAG (bilgi) + hafif fine-tune (davranış)} birlikte kullanır.

\stage{11}{İlişkili Kavramlar}
\term{RAG} (bilgi tarafı) \tlsep \term{Prompt} (önce tüketilecek) \tlsep
\term{LoRA / PEFT} (ucuz fine-tune) \tlsep \term{Distillation} (büyük modeli küçüğe öğretme) \tlsep
\term{Eval} (fine-tune gerçekten iyileştirdi mi?)

\stage{12}{Özet}
\begin{itemize}
  \item Fine-tuning \emph{davranışı} öğretir, \emph{bilgiyi} değil.
  \item Bilgi $\to$ RAG; stil/format/ton $\to$ fine-tune. Karıştırma.
  \item Önce prompt'u tüket; fine-tune en pahalı ve en kalıcı seçenektir.
  \item Çıktı bir model ağırlığıdır: bilgi değişince yeniden eğitim gerekir.
  \item En sağlam kombinasyon: RAG + hafif fine-tune.
\end{itemize}

\stage{13}{Mühendislik İlkesi}
\begin{principlebox}
\textbf{İlke \#5 — Fine-tuning davranış içindir, bilgi için değil.}\\[4pt]
\normalsize Modele ne \emph{bildiğini} değil, nasıl \emph{davrandığını} öğret.
Bilgi taze tutulmak ister; onu ağırlıklara değil, retrieval'a emanet et.
\end{principlebox}
